\section{A crossing-rank formulation of the square question}
\label{sec:crossing-rank}

This section puts the square question into a classical algebraic
form: the rank of an identity matrix after decomposing it across
two different cuts. The relevant notion is the partial-transpose
rank of Rossman and Zhu \cite[Definition 3.4 and Lemma 3.5]{RossmanZhu2026}.
The upper exponent $\sqrt2$ is already obtained from crossing
identities by Khoruzhii \cite[Lemma 1 and Corollary 3]{Khoruzhii2026}.
Liu proves that the absence of a planar strict shortening forces
the complex Hurwitz exponent to equal $\sqrt2$, using the two
adjacent cuts and a two-hub routing argument
\cite[Section 2, p.~3, ``Two crossing cuts'']{Liu2026noncommutative}.
Here this sharpness statement is related to the square-price
question, with a converse and uniform eventual bounds over every
field, assuming the price-space theorem.

For an integer $q\geq2$ put
\[
 C_q=\sum_{a,b=1}^q e_a\otimes e_b\otimes e_a\otimes e_b.
\]
For a four-party tensor $T$, write $r_0(T)$ for its matrix rank
across $12\mid34$ and $r_1(T)$ for its rank across $14\mid23$.
Define
\begin{equation}\label{eq:pt-crossing-definition}
 P_\kk(q)=\min_{C_q=A+B}\bigl(r_0(A)+r_1(B)\bigr).
\end{equation}
Zero summands are allowed. On identifying the $12\mid34$
flattening with $q^2$ by $q^2$ matrices, $C_q$ is the identity
matrix and changing the second cut is partial transposition.
Thus \eqref{eq:pt-crossing-definition} is precisely its
two-factor partial-transpose rank. Indeed a rank decomposition
of $A$ and of the second flattening of $B$ gives that many
partial-transpose-basic summands; conversely group all such
summands by their two possible cuts.

Naslund's Problem~12 asks for a superlinear lower bound on this
same quantity \cite{Naslund2020}. His convention uses the tensor
$\delta_{xv}\delta_{yu}$ and cuts $xy\mid uv$ and $xu\mid yv$;
the substitution $(a,b,c,d)=(x,y,v,u)$ gives
\eqref{eq:pt-crossing-definition}.
Singh obtains the all-field bounds $P_\kk(q)\geq4q/3$ for $q\geq3$,
$P_\kk(q)\geq67q/49$ for $q\geq7$, and the exact value
$P_\kk(4)=7$ \cite[Theorems~1--4 and Corollary~5, p.~2]{Singh2026crossing}.
These are linear or finite-dimensional results and do not determine
the exponent.\footnote{The complete subspace proof of $4q/3$ in
\cite[Section~3, pp.~3--5]{Singh2026crossing} was checked directly.
The source provides exact certificates for the larger linear constants
and Lean proofs for its exact small-dimensional assertions; those
Lean proofs have not been replayed here.}

\begin{lemma}[The known crossing upper bound]\label{lem:pt-upper}
Over every field,
\[
 q\leq P_\kk(q)\leq q^{\sqrt2+o(1)}\qquad(q\longrightarrow\infty).
\]
\end{lemma}
\begin{proof}
Contract the first and third parties of $C_q$ against the same
coordinate vector. The result is the identity matrix of order
$q$ on the other two parties. The corresponding contraction
of $A$ has rank at most $r_0(A)$, and that of $B$ has rank at
most $r_1(B)$. Rank subadditivity proves the lower bound.

For clarity we give the counting passage from the integral
crossing identity in \cite[Proposition B.2]{Anthropic2026}.
It is the same exponent calculation as in
\cite[Lemma 1 and Corollary 3]{Khoruzhii2026}.
Among the symbol classes of energy at most $E$, choose one of
largest cardinality $q_E$. There are polynomially many classes,
and their total count has logarithm $(\kappa+o(1))\sqrt E$.
Consequently
\[
 \log q_E=(\kappa+o(1))\sqrt E.
\]
Use that class for both crossing wires. If its energy is $e\leq E$,
the identity has $4e+1$ coefficient tensors. Put in $A$ the
coefficients of index at most $2e$ and in $B$ the others,
using the corresponding adjacent cut for each group.
Each required cut rank is at most
$(4e+1)^2P(2e)$, with the symbol-counting function $P$
from Section~\ref{sec:curvature}. Therefore
\[
 P_\kk(q_E)\leq 2(4E+1)^2P(2E)
             =\exp\bigl((\kappa\sqrt2+o(1))\sqrt E\bigr).
\]
This argument is integral and hence valid in every characteristic.
To pass to every $q$, fix $\delta>0$ and take
$E=\lceil((1+\delta)\log q/\kappa)^2\rceil$.
For large $q$ the chosen class has $q_E\geq q$.
Restricting all four index sets to their first $q$ elements
cannot increase either cut rank. The resulting upper exponent
is at most $(1+\delta)\sqrt2+o(1)$.
Letting $\delta$ decrease to zero proves the assertion.
\end{proof}

The sharp-exponent direction below includes the routing argument
of \cite[Section 2, p.~3]{Liu2026noncommutative}; the price-space
duality also supplies the converse and the uniform improvement.

\begin{theorem}[The square and the sharp crossing exponent]
\label{thm:square-pt-exponent}
Fix a field $\kk$ and assume the price-space theorem
\cite[Proposition C.3]{Anthropic2026}.
The unit square is a price over $\kk$ if and only if
\begin{equation}\label{eq:pt-sharp}
 \lim_{q\to\infty}\frac{\log P_\kk(q)}{\log q}=\sqrt2.
\end{equation}
If the square is not a price, there is a $\delta>0$ for which
\begin{equation}\label{eq:pt-power-improvement}
 P_\kk(q)\leq q^{\sqrt2-\delta}
 \quad\hbox{for all sufficiently large }q.
\end{equation}
In particular, an improvement below $\sqrt2$ along any unbounded
sequence of dimensions by a fixed positive amount already
implies such a uniform eventual improvement.
\end{theorem}
\begin{proof}
Suppose first that the square is not a price. The proof of
Theorem~\ref{thm:square-shortening-equivalence} gives positive
integers $m,h$ with $h/m<2\sqrt2$ such that, for every
$\varepsilon>0$ and all large $N$, $C_{2^{mN}}$ is covered
by at most $2^{\varepsilon N}$ restrictions of side-wire
products of total width $hN$.
For one such product, the widths crossing the two adjacent
cuts sum to $hN$: the four sides are partitioned between
those cuts. One cut therefore has width at most $hN/2$.
The corresponding restriction has cut rank at most
$2^{hN/2}$. Grouping the cover terms according to their
cheaper cut gives
\[
 P_\kk(2^{mN})\leq 2^{(h/2+\varepsilon)N}.
\]
Choose $\varepsilon>0$ so that
$\alpha=(h/2+\varepsilon)/m<\sqrt2$.
For arbitrary large $q$, choose
$N=\lceil\log_2q/m\rceil$ and restrict from $2^{mN}$ to $q$.
Since $mN/\log_2q\to1$, this proves
\eqref{eq:pt-power-improvement} with, for example, any
$0<\delta<\sqrt2-\alpha$.

Conversely, suppose that for some $\alpha<\sqrt2$ there are
arbitrarily large $q$ with $P_\kk(q)\leq q^\alpha$.
Choose a minimizing decomposition in
\eqref{eq:pt-crossing-definition}, and put
$m=\lfloor\log_2q\rfloor$. Restrict all parties to dimension
$2^m$. Both summands have their indicated cut rank at most
$q^\alpha$, so each can be routed through the corresponding
two adjacent-pair hubs with bridge width
\[
 h\leq\lceil\alpha\log_2q\rceil=\alpha m+O(1).
\]
All four access cables have width $m$.
Proposition~\ref{prop:optimal-crossing-hubs} places the hubs
in the unit disc and bounds each alternative cost by
$mF(h/m)$. Choose $\beta$ with
$\alpha<\beta<\sqrt2$. For all large chosen $q$ we have
$h/m\leq\beta$ and $F(h/m)\leq F(\beta)<4$.
There are at most two alternatives, whose centre selector
charge is at most $4$. Hence
\[
 mF(h/m)+4<4m
\]
for sufficiently large $q$. This is a strict shortening,
so the square is not a price by
Theorem~\ref{thm:square-shortening-equivalence}.

If the square is a price, the second argument shows that
the lower limit in \eqref{eq:pt-sharp} is at least $\sqrt2$.
Lemma~\ref{lem:pt-upper} gives the reverse bound for the upper
limit. If the square is not a price, the first argument
precludes \eqref{eq:pt-sharp}. These prove all assertions.
No submultiplicativity of $P_\kk(q)$ under grouped tensor
powers is used.
\end{proof}

\subsection{Hurwitz's problem over the complex numbers}

Let $\sigma_\C(q)$ be the least number of complex bilinear forms
$f_\nu(x,y)$ in two blocks of $q$ variables such that
\[
 \left(\sum_{i=1}^q x_i^2\right)
 \left(\sum_{j=1}^q y_j^2\right)
      =\sum_{\nu=1}^{\sigma_\C(q)}f_\nu(x,y)^2.
\]
This is Hurwitz's sum-of-squares composition problem, in its
complex form. Hrube\v{s}'s subquadratic construction
\cite[Theorem 1]{Hrubes2026} and the subsequent crossing bound
\cite[Corollary 3]{Khoruzhii2026} concern exactly this quantity.
The comparison with partial-transpose rank is developed in
\cite[Lemmas 3.12--3.14]{RossmanZhu2026}.

\begin{corollary}[The square and the Hurwitz exponent]
\label{cor:square-hurwitz}
Assuming the complex price-space theorem, the square and the
square with its centre are prices if and only if
\[
 \lim_{q\to\infty}\frac{\log\sigma_\C(q)}{\log q}=\sqrt2.
\]
If they are not prices, then
$\sigma_\C(q)\leq q^{\sqrt2-\delta}$ for all sufficiently
large $q$, for some $\delta>0$.
\end{corollary}
\begin{proof}
We prove an elementary constant-factor comparison sufficient
for the exponent.
Each rank-one term across either adjacent cut becomes a product
$\ell(x,y)m(x,y)$ of two bilinear forms when the first and
third variable blocks, and the second and fourth blocks, are
identified. Over $\C$,
\[
 \ell m=\left(\frac{\ell+m}{2}\right)^2+
         \left(\frac{\ell-m}{2i}\right)^2.
\]
Thus $\sigma_\C(q)\leq2P_\C(q)$.

Conversely, a formula with $s$ squares has a symmetric Gram
matrix $G=\sum_{\nu=1}^s v_\nu v_\nu^{\mathsf T}$ of rank
at most $s$, where $v_\nu$ lists the bilinear coefficients.
Polarizing separately in the two variable blocks gives
\[
 I_{q^2}=\tfrac12(G+G^\Gamma),
\]
where $\Gamma$ is partial transposition. To check the factor,
the full polarization is the average of the four partial
transposes; $G$ and $G^\Gamma$ are symmetric, so the four
terms consist of two copies of each. The two displayed
summands have respective cut ranks at most $s$, giving
$P_\C(q)\leq2\sigma_\C(q)$.
Multiplicative constants disappear after taking logarithms
and dividing by $\log q$. Apply
Theorem~\ref{thm:square-pt-exponent} and the planar equivalence.
\end{proof}

This reformulation must not be transferred to characteristic
two. There the polynomial in Hurwitz's problem is itself the
single square $(\sum_i x_i\sum_j y_j)^2$, while
$P_\kk(q)\geq q$ still holds. Polarization loses the crossing
tensor in that characteristic. The partial-transpose formulation
\eqref{eq:pt-crossing-definition} remains valid over every field.

\subsection{Why two side-wire alternatives cannot suffice}

The normal form in Theorem~\ref{thm:square-shortening-equivalence}
is complete for detecting a nonprice square. The next obstruction
rules out its smallest nontrivial number of alternatives.

\begin{proposition}[Two alternatives in the side-wire normal form]
\label{prop:two-side-alternatives}
Let $m\geq1$ be an integer. Suppose $C_{2^m}$ is covered over a field by at
most two restrictions of products of wires on the four sides
of the square, with every wire of dimension $2^{b_e}$ for a
nonnegative integer $b_e$. If every product has total width at most $b$,
then, in the two-term case, $b\geq3m-1$.
No such one- or two-term cover is a strict shortening at the
unit-circle terminal positions. The conclusion also holds for
the Zariski closure of the covers with these width bounds.
\end{proposition}
\begin{proof}
We may replace $b$ by the maximum actual total width, which is
an integer. Put $q=2^m$. For alternative $j$ let $s_j,t_j$ be the total
widths across $12\mid34$ and $14\mid23$. Every side crosses
exactly one of these cuts, so $s_j+t_j\leq b$.
Set $u_j=2^{s_j}$ and $v_j=2^{t_j}$.
The crossing tensor has rank $q^2$ across each cut, whence
\[
 q^2\leq u_1+u_2,\qquad q^2\leq v_1+v_2.
\]
Contract parties $1$ and $3$ against the same coordinate
vector. The crossing tensor becomes a rank-$q$ matrix on
parties $2,4$. Each alternative's contracted matrix has rank
at most each of its two cut ranks. Hence
\[
 q\leq\min(u_1,v_1)+\min(u_2,v_2).
\]
The first two bounds force $s_j\geq2m-1$ for some $j$
and $t_k\geq2m-1$ for some $k$.
Suppose, for a contradiction, that $b\leq3m-2$.
These two large widths cannot occur on the same alternative,
since $4m-2>3m-2$. Relabel so that
$s_1,t_2\geq2m-1$, and hence $t_1,s_2\leq m-1$.
If $s_1=2m-1$, then
\[
 2^{s_1}+2^{s_2}\leq2^{2m-1}+2^{m-1}<2^{2m},
\]
contrary to the first cut-rank bound.
Thus $s_1\geq2m$, and similarly $t_2\geq2m$.
It follows that $t_1,s_2\leq m-2$.
For $m=1$ this contradicts nonnegative widths.
For $m\geq2$ the contracted-wire bound would give
$2^m\leq2^{m-2}+2^{m-2}=2^{m-1}$, also impossible.
Therefore $b\geq3m-1$.
At $m=1$ the two-term Pl\"ucker cover attains $b=2$.
The geometric cost including the two-term selector is at least
\[
 \sqrt2(3m-1)+4
 =4m+(3\sqrt2-4)m+4-\sqrt2>4m.
\]
For one alternative, both cut ranks require
$s_1,t_1\geq2m$, so $b\geq4m$, which also precludes shortening.

For fixed widths the three rank bounds used here are closed
conditions given by minors, including the minors after the
fixed contraction. They therefore hold on the Zariski closure
of the corresponding set of covers. Only finitely many
integral width choices satisfy the bound $b$; closure commutes
with a finite union. The same argument applies to that union.
\end{proof}
